

\section{Les Seuils de Lecture Fluide : Objectif 95-98% de mots connus}
space{0.3cm}
oindent

La crise contemporaine de l'enseignement des langues anciennes, et singulièrement du grec ancien,
trouve l'une de ses explications majeures dans la persistance d'un modèle pédagogique centré sur le
déchiffrage intensif (intensive reading) au détriment de la fluidité de lecture (extensive reading).
Le projet de refondation pédagogique, dont le point III.1.\cite{III-1-2-ref2} postule un \og
Objectif 95-98\% de mots connus pour lire sans dictionnaire\fg{}, marque une rupture épistémologique
nécessaire. Il ne s'agit plus de considérer le texte grec comme un cryptogramme grammatical à
décoder, mais comme un flux d'informations linguistiques à traiter en temps réel. Cette transition
exige une compréhension fine des mécanismes cognitifs sous-jacents à la lecture et une
quantification rigoureuse des seuils lexicaux qui permettent l'autonomie.

Ce rapport se propose d'analyser de manière exhaustive les fondements scientifiques de ces seuils de
95\% et 98\%, en s'appuyant sur la littérature en acquisition des langues secondes (SLA), la
psycholinguistique cognitive et la statistique lexicale appliquée aux corpus grecs (Attique, Koinè,
et corpus bibliques). L'enjeu est de démontrer que la lecture sans dictionnaire n'est pas une simple
commodité, mais une condition \textit{sine qua non} pour l'activation des processus d'inférence
contextuelle et l'acquisition implicite de la syntaxe. Nous examinerons pourquoi le seuil de 98\%
constitue l'optimum pour une lecture de plaisir et d'apprentissage durable, tandis que le seuil de
95\% représente un minimum fonctionnel souvent insuffisant pour les textes à haute densité
conceptuelle comme ceux de Platon ou de Thucydide.\cite{III-1-2-ref1}

L'analyse portera également sur les obstacles spécifiques posés par la morphologie du grec ancien,
langue flexionnelle où l'opacité des formes peut masquer la transparence lexicale, imposant des
contraintes supplémentaires par rapport aux langues moins synthétiques. Enfin, nous évaluerons
l'adéquation des ressources pédagogiques actuelles face à ces exigences statistiques, mettant en
lumière le \og fossé du niveau intermédiaire\fg{} qui condamne trop souvent les apprenants à une
dépendance perpétuelle aux outils de référence.\cite{III-1-2-ref4}



\subsection{Fondements théoriques : la psycholinguistique des seuils lexicaux}

La détermination des seuils de couverture lexicale nécessaires à la compréhension n'est pas
arbitraire. Elle résulte de décennies de recherches empiriques visant à quantifier la relation entre
la densité des mots inconnus et la qualité de la compréhension textuelle. Cette section détaille la
distinction cruciale entre le seuil de \og compréhension adéquate\fg{} et le seuil de \og lecture
fluide\fg{}.



\subsection{Le seuil minimal de 95\% : la frontière de la compréhension assistée}

Les travaux fondateurs de Batia Laufer (1989) ont établi un premier consensus autour du seuil de
95\%. Selon ces recherches, un lecteur doit connaître au moins 95\% des unités lexicales (tokens)
d'un texte pour parvenir à une compréhension globale satisfaisante.\cite{III-1-2-ref6} Concrètement,
cela signifie que le lecteur rencontre un mot inconnu tous les 20 mots, soit environ un mot par
phrase ou toutes les deux lignes dans un texte grec standard.

À ce niveau de couverture (95\%), la compréhension est possible mais elle demeure fragile. Le
lecteur est capable de saisir la macro-structure du texte et les idées principales, mais il doit
engager des ressources cognitives importantes pour compenser les lacunes lexicales. Les études
montrent que l'inférence contextuelle (la capacité de deviner le sens d'un mot grâce à son
environnement) est activée, mais son taux de succès est variable et dépend fortement de la
transparence du contexte immédiat.\cite{III-1-2-ref8} Pour un apprenant du grec ancien, souvent
confronté à des structures syntaxiques complexes (hyperbates, participes enchâssés), un mot inconnu
tous les 20 mots peut suffire à obscurcir les relations grammaticales essentielles, transformant la
lecture en une tâche de résolution de problèmes plutôt qu'en une réception fluide du
sens.\cite{III-1-2-ref9}

Il est donc impératif de considérer le seuil de 95\% comme un seuil \textit{minimal} ou
\textit{instructif}, adapté à une lecture accompagnée en classe ou à l'étude intensive, mais
insuffisant pour garantir l'autonomie complète de l'élève face à un texte authentique
complexe.\cite{III-1-2-ref2}



\subsection{Le seuil optimal de 98\% : la condition de la fluidité et du plaisir}

La recherche a connu un tournant décisif avec l'étude de Hu et Nation (2000), intitulée
\textit{Unknown Vocabulary Density and Reading Comprehension}. En utilisant des textes où les mots
étaient remplacés par des non-mots à différentes densités, ils ont démontré que pour une lecture non
assistée (unassisted reading), fluide et permettant le plaisir (reading for pleasure), le seuil de
couverture doit être porté à \textbf{98\%}.\cite{III-1-2-ref1}

Ce seuil de 98\% implique que le lecteur ne rencontre qu'un seul mot inconnu tous les 50 mots. Cette
différence de trois points de pourcentage (de 95\% à 98\%) change radicalement l'expérience
cognitive de la lecture :

\begin{itemize}\item \textbf{Automaticité de l'inférence :} À 98\%, les indices contextuels sont suffisamment abondants et non ambigus pour que le cerveau puisse prédire ou inférer le sens du mot manquant avec une haute probabilité de succès, sans rupture significative du flux attentionnel.\cite{III-1-2-ref3}\item \textbf{Réserve attentionnelle :} Le lecteur dispose de suffisamment de ressources en mémoire de travail pour traiter non seulement le sens littéral, mais aussi les nuances stylistiques, rhétoriques et philosophiques du texte, dimension essentielle pour la littérature grecque antique.\cite{III-1-2-ref2}\item \textbf{Endurance :} Contrairement à la lecture à 95\% qui est fatigante, la lecture à 98\% peut être soutenue sur de longues périodes. C'est cette endurance qui permet l'exposition massive (input flood) nécessaire pour que les structures grammaticales complexes du grec (comme le système des temps et des aspects) soient intériorisées par imprégnation plutôt que par analyse consciente.\cite{III-1-2-ref12}\end{itemize}Laufer et Ravenhorst-Kalovski (2010) ont confirmé cette distinction en proposant deux seuils
distincts : 95\% pour une compréhension minimale et 98\% pour une compréhension
optimale.\cite{III-1-2-ref2} Dans le contexte d'une refondation de l'enseignement du grec, viser le
98\% pour les textes de lecture cursive est le seul moyen de sortir de la logique de traduction mot
à mot.



\subsection{Nuances méthodologiques et validité pour les langues anciennes}

Il convient de noter que la majorité de ces études ont été menées sur l'anglais langue seconde.
Cependant, les principes cognitifs sous-jacents (limites de la mémoire de travail, traitement de
l'information) sont universels. Pour le grec ancien, la situation pourrait même exiger une rigueur
accrue. Van Zeeland et Schmitt (2013) notent que si 95\% peut suffire pour une compréhension
auditive grâce aux indices prosodiques, la lecture de textes denses et distants culturellement ne
bénéficie pas de ces aides paralinguistiques.\cite{III-1-2-ref7}

De plus, la nature de l'évaluation de la compréhension influence les résultats. Hu et Nation ont
défini une compréhension adéquate comme l'obtention de 12/14 à un test à choix multiples (QCM) ou
70/124 à un test de rappel écrit.\cite{III-1-2-ref12} Ces standards élevés confirment que pour lire
Platon ou le Nouveau Testament avec la précision requise par l'exégèse ou l'analyse philosophique,
le seuil de 98\% est un objectif technique incontournable, et non une simple préférence pédagogique.



\subsection{La psychologie cognitive : charge cognitive et dictionnaire}

L'injonction de \og lire sans dictionnaire\fg{} ne relève pas d'un purisme académique, mais d'une
nécessité cognitive. La Théorie de la Charge Cognitive (Cognitive Load Theory \- CLT) fournit un
cadre explicatif puissant pour comprendre pourquoi la consultation externe d'outils lexicaux entrave
l'acquisition de la fluidité.



\subsection{Le coût cognitif de la consultation lexicale}

L'acte de consulter un dictionnaire, qu'il soit papier ou numérique, impose une \og double
tâche\fg{} au cerveau de l'apprenant. Lorsqu'un lecteur interrompt sa lecture pour chercher un mot :

1. \textbf{Rupture de la Boucle Phonologique :} La phrase grecque maintenue en mémoire de travail
(boucle phonologique) commence à s'estomper. La capacité de la mémoire de travail étant limitée
(généralement estimée à quelques éléments pendant quelques secondes), le temps passé à chercher le
mot dépasse souvent la durée de rétention de la structure syntaxique en cours de
traitement.\cite{III-1-2-ref15}

2. \textbf{Surcharge Cognitive Extrinsèque :} La recherche alphabétique, le tri des différentes
acceptions d'un mot polysémique (très fréquent en grec, cf. \textit{logos}, \textit{archè}), et la
sélection du sens approprié constituent une charge cognitive \og extrinsèque\fg{} qui n'est pas
directement liée à la compréhension du texte mais à la gestion de l'outil.\cite{III-1-2-ref17}

3. \textbf{Effet de Redémarrage :} Une fois le sens trouvé, l'élève doit souvent relire la phrase
depuis le début pour réintégrer le mot dans son contexte syntaxique, doublant ainsi le temps de
traitement et brisant le rythme nécessaire à la compréhension globale.\cite{III-1-2-ref18}

Les études sur l'utilisation des dictionnaires montrent que, bien que la consultation puisse mener à
une meilleure rétention du mot spécifique cherché (grâce à l'effort de traitement, ou \og depth of
processing\fg{}), elle nuit significativement à la fluidité de lecture et à la compréhension des
relations logiques du texte.\cite{III-1-2-ref19} Prichard (2008) souligne que pour la lecture de
plaisir ou le développement de la fluidité, les lecteurs ne devraient idéalement pas avoir à
consulter de dictionnaire du tout.\cite{III-1-2-ref13}



\subsection{L'effet d'attention divisée (split-attention effect)}

La présentation matérielle du vocabulaire joue un rôle crucial. La CLT identifie l'effet d'attention
divisée (\textit{split-attention effect}) comme une source majeure de difficulté d'apprentissage.
Lorsque le vocabulaire est présenté sous forme de listes séparées (en fin de volume ou dans un
fascicule distinct, comme c'était le cas pour les premières éditions de \textit{Reading Greek}),
l'apprenant doit constamment alterner son attention visuelle entre le texte et la
liste.\cite{III-1-2-ref21} Cette intégration mentale forcée consomme des ressources cognitives
précieuses.

Les recherches en design instructionnel démontrent que l'intégration physique des gloses (par
exemple, en marge, en bas de page immédiat, ou via des hyperliens contextuels) réduit cette charge
et favorise une meilleure compréhension.\cite{III-1-2-ref17} Cependant, l'objectif du seuil de 98\%
est de dépasser même le besoin de ces gloses intégrées, pour que l'accès au sens soit direct et
automatisé. L'étude de Yeung (1997) confirme que si les formats intégrés sont supérieurs aux formats
divisés pour les novices, l'objectif ultime reste l'internalisation du lexique pour libérer la
mémoire de travail au profit des processus de haut niveau (analyse littéraire, critique,
théologique).\cite{III-1-2-ref17}



\subsection{La réalité statistique du grec ancien : le mur de zipf et les hapax}

Pour opérationnaliser les seuils de 95-98\% dans une refondation curriculaire, il est indispensable
de confronter ces pourcentages à la réalité statistique des corpus grecs. La distribution du
vocabulaire en grec ancien suit, de manière impitoyable, la loi de Zipf, créant des défis
spécifiques pour l'atteinte des seuils de fluidité.



\subsection{La loi de zipf et la longue traîne des mots rares}

La loi de Zipf stipule que la fréquence d'un mot est inversement proportionnelle à son rang dans la
table des fréquences. Concrètement, cela signifie qu'un petit noyau de mots très fréquents
(articles, pronoms, particules, verbes être/dire/faire) couvre une proportion massive du texte,
tandis qu'une immense quantité de mots rares (\textit{hapax legomena} et \textit{dis legomena})
constitue le reste.\cite{III-1-2-ref24}

\begin{itemize}\item \textbf{Les Mots de Haute Fréquence :} Les \~500 mots les plus fréquents du grec couvrent environ 65\% à 70\% de n'importe quel texte standard.\cite{III-1-2-ref27} L'apprentissage de ce noyau est rapide et gratifiant.\item \textbf{La Zone de Rendement Décroissant :} Au-delà de ces 500 mots, chaque nouveau mot appris apporte une contribution marginale de plus en plus faible à la couverture totale. C'est ici que se situe le piège pédagogique : l'étudiant a l'impression de connaître \og beaucoup de mots\fg{} mais continue de buter sur des inconnus à chaque phrase.\cite{III-1-2-ref29}\end{itemize}

\subsection{Quantification précise des seuils (données tauber \& macdonald)}

Les travaux récents de James Tauber et Seumas Macdonald, basés sur l'analyse computationnelle des
corpus (Projet MorphGNT, Corpus Platonicien), fournissent les métriques exactes nécessaires pour
calibrer le programme d'enseignement.

Le tableau ci-dessous synthétise le nombre de \textbf{lemmes} (mots du dictionnaire) nécessaires
pour atteindre les seuils critiques de 95\% et 98\% dans différents corpus :

Tableau 1 : Nombre de lemmes requis pour les seuils de couverture 4

\begin{table}[h!]\small\centering\begin{tabular}{| p{2.3cm} | p{3.5cm} | p{3.5cm} |}\hline\textbf{Corpus Cible} & \textbf{Seuil 95\% (Compréhension minimale)} & \textbf{Seuil 98\% (Lecture Fluide)} \\ \hline\textbf{Nouveau Testament (GNT)} & \~1 753 lemmes & \textbf{\~3 103 lemmes} \\ \hline\textit{Platon (Sélection)}\* & \~1 745 lemmes & \textbf{\~3 160 lemmes} \\ \hline\textbf{Platon (Corpus général)} & \~1 840 lemmes & \textbf{\~3 631 lemmes} \\ \hline\end{tabular}\end{table}\textit{\}Sélection : Euthyphron, Apologie, Criton, Phédon, République.*

1. \textbf{Le Plateau Intermédiaire :} La plupart des manuels introductifs traditionnels (comme
\textit{Reading Greek} ou \textit{Athenaze} vol. 1) enseignent entre 1 000 et 1 500
mots.\cite{III-1-2-ref22} Cela place l'étudiant en dessous du seuil de 95\% (autour de 85-90\%), le
laissant dans une zone de frustration cognitive intense où la lecture autonome est impossible.

2. \textbf{Le Coût du Passage à 98\% :} Pour passer d'une compréhension assistée (95\%) à une
lecture fluide (98\%), l'étudiant doit presque \textbf{doubler} son vocabulaire, passant de \~1 750
à \~3 100 mots. Ce \og kilomètre supplémentaire\fg{} est souvent négligé dans les cursus, créant un
plafond de verre.

3. \textbf{Convergence des Corpus :} Il est remarquable de noter la similarité des chiffres entre le
Nouveau Testament et Platon (\~3 100 lemmes pour 98\%). Cela suggère qu'un \og vocabulaire
noyau\fg{} de 3 000 à 3 500 mots constitue un objectif curriculaire robuste pour une compétence de
lecture polyvalente en grec ancien (Attique et Koinè).\cite{III-1-2-ref4}



\subsection{Le défi morphologique : au-delà du lemme}

L'analyse purement lexicale masque une difficulté intrinsèque au grec ancien : sa morphologie
flexionnelle complexe. Dans les langues comme l'anglais ou l'espagnol, connaître le mot de base
suffit souvent à reconnaître ses formes dérivées. En grec, la distance orthographique et
phonologique entre le lemme et la forme fléchie peut être immense, compromettant la reconnaissance
immédiate nécessaire à la fluidité.



\subsection{L'opacité morphologique comme obstacle à la reconnaissance}

La recherche psycholinguistique sur les langues flexionnelles (comme le finnois ou le grec) montre
que la reconnaissance des mots ne procède pas toujours par décomposition analytique (parsing) chez
les apprenants L2, mais souvent par reconnaissance globale (holistique).\cite{III-1-2-ref31}

\begin{itemize}\item \textbf{Allomorphie et Irrégularité :} Un étudiant peut connaître le verbe \textit{ferō} (porter) et \textit{horaō} (voir), mais échouer à identifier \textit{enēnegmai} (parfait passif de \textit{ferō}) ou \textit{ōphthēn} (aoriste passif de \textit{horaō}) comme appartenant à ces lemmes. Pour le système cognitif de l'apprenant, ces formes fonctionnent comme des mots distincts, augmentant la charge mémorielle réelle bien au-delà des 3 100 lemmes théoriques.\cite{III-1-2-ref5}\item \textbf{Traitement des Formes Opaques :} Les études indiquent que les apprenants intermédiaires ont des temps de fixation plus longs et des taux d'erreur plus élevés sur les formes fléchies opaques (fusionnelles) par rapport aux formes agglutinantes transparentes.\cite{III-1-2-ref31} Cela ralentit la lecture et perturbe l'accès au sens, même si le \og vocabulaire\fg{} est théoriquement acquis.\end{itemize}

\subsection{La conscience morphologique (morphological awareness)}

Pour atteindre le seuil de 98\% de \og mots connus\fg{} en situation réelle de lecture, la
refondation pédagogique ne peut se contenter d'enseigner des listes de lemmes. Elle doit développer
explicitement la \textbf{conscience morphologique} (morphological awareness).\cite{III-1-2-ref5}

\begin{itemize}\item \textbf{Enseignement Explicite du Parsing :} Les données suggèrent que la capacité à décomposer un mot inconnu en morphèmes (racine \+ suffixes dérivationnels \+ désinences flexionnelles) est un prédicteur fort de la compréhension de lecture, parfois supérieur à la simple étendue du vocabulaire chez les débutants.\cite{III-1-2-ref5}\item \textbf{Stratégie d'Intervention :} L'enseignement doit se focaliser sur les \og systèmes de racines\fg{} et les règles de formation des mots (dérivation par préfixes et suffixes). Par exemple, reconnaître que les composés en \textit{a-} privatif, \textit{eu-}, \textit{dys-}, ou les substantifs en \textit{\-sis}, \textit{\-ma}, \textit{\-tēs} sont déductibles, permet de transformer des \og mots inconnus\fg{} lexicaux en \og mots connus\fg{} morphologiquement, augmentant artificiellement le taux de couverture du lecteur.\cite{III-1-2-ref34}\end{itemize}La maîtrise des 3 100 lemmes doit donc s'accompagner d'une maîtrise automatisée des paradigmes
verbaux et nominaux pour que le seuil de 98\% soit effectif sur le texte \textit{tel qu'il apparaît}
(surface forms) et non seulement dans l'abstrait du dictionnaire.



\subsection{État des lieux critique des ressources pédagogiques}

Si l'objectif est d'amener les étudiants à un seuil de 3 000+ mots actifs pour une lecture fluide,
les outils actuels sont-ils adaptés? L'analyse comparative des manuels révèle des disparités
majeures dans leur capacité à soutenir cette ambition.



\subsection{Les limites des manuels "reading method" classiques}

\begin{itemize}\item \textbf{JACT \textit{Reading Greek} :} Pionnier de la méthode de lecture, ce manuel a le mérite d'introduire des textes continus dès le début. Cependant, son vocabulaire total (environ 1 300 à 1 500 mots) est insuffisant pour le saut vers les textes authentiques non adaptés.\cite{III-1-2-ref36} De plus, la séparation initiale du texte et du vocabulaire (dans l'édition originale) induisait un fort effet d'attention divisée.\cite{III-1-2-ref38} Les critiques soulignent une courbe de progression trop abrupte, où la densité de mots inconnus dépasse rapidement les seuils de tolérance cognitive pour l'étude autonome.\cite{III-1-2-ref38}\item \textbf{Athenaze (Édition US/UK) :} Bien que structuré autour d'une narration, \textit{Athenaze} standard souffre d'un volume de texte insuffisant pour consolider le vocabulaire par la répétition (input flood). Le vocabulaire cible est également limité, laissant l'étudiant sur le \og plateau intermédiaire\fg{}.\cite{III-1-2-ref40}\end{itemize}

\subsection{Les nouvelles approches : vers l'extensif et l'immersif}

\begin{itemize}\item \textbf{Athenaze (Édition Italienne \- Miraglia) :} Cette adaptation représente une avancée majeure. En doublant le volume de texte (près de 4 000 lignes dans le vol.\cite{III-1-2-ref1} contre \~1 500 dans l'édition US) et en intégrant un vocabulaire plus vaste (\~2 600 mots à la fin du vol. 2), elle rapproche considérablement l'étudiant du seuil d'autonomie pour le Nouveau Testament et la prose simple.\cite{III-1-2-ref30} L'approche \og Orbergienne\fg{} (définitions en grec ou illustrées en marge) réduit la charge cognitive extrinsèque et favorise l'immersion.\item \textbf{Alexandros : To Hellenikon Paidion :} Basé sur l'ouvrage de Rouse (\textit{A Greek Boy at Home}), ce lecteur offre une expérience de lecture extensive précieuse (entre 5 000 et 13 000 mots de texte).\cite{III-1-2-ref43} Il permet de saturer le vocabulaire de base dans des contextes variés. Cependant, le vocabulaire de Rouse est parfois idiosyncratique et ne correspond pas parfaitement aux fréquences statistiques des corpus classiques cibles, nécessitant une médiation de l'enseignant.\cite{III-1-2-ref45}\item \textbf{Méthode Polis (Rico) :} Axée sur l'oralité et la fluidité (Fluency), cette méthode vise explicitement l'internalisation des structures. Bien qu'efficace pour l'automatisation, elle nécessite des compléments de lecture extensive pour atteindre le volume lexical de 3 000 lemmes.\cite{III-1-2-ref46}\end{itemize}

\subsection{Le projet "greek learner texts" et le manque de "ponts"}

Le constat global est celui d'un manque de textes \og passerelles\fg{} (bridge texts). Il existe des
ressources pour débutants (0-1000 mots) et des textes authentiques (\>3000 mots), mais très peu de
matériel calibré pour la zone 1000-3000 mots, essentielle pour atteindre les 98\%.

Le projet Greek Learner Texts (Tauber, Macdonald) tente de combler ce vide en produisant des corpus
annotés et segmentés, mais le volume disponible reste en deçà des besoins pour une véritable \og
lecture extensive\fg{} qui exigerait des centaines de pages à ce niveau
intermédiaire.\cite{III-1-2-ref48} Macdonald souligne l'impossibilité de faire lire 280 minutes par
semaine à des débutants sans textes à \og vocabulaire abrité\fg{} (sheltered
vocabulary).\cite{III-1-2-ref50}



\subsection{Stratégies de refondation : l'ingénierie curriculaire du 98\%}

Pour que l'objectif III.1.\cite{III-1-2-ref2} devienne réalité, la refondation doit dépasser la
simple sélection de manuels pour adopter une ingénierie curriculaire basée sur les données.



\subsection{Calibrage statistique du vocabulaire}

Le programme doit être piloté par les données de fréquence. Il convient d'adopter les listes de
fréquence curées (comme celles de Tauber ou du DCC Core Vocabulary) pour définir les étapes
d'apprentissage :

1. \textbf{Niveau 1 (Fondation) :} Les 500 mots les plus fréquents (couverture
\~65-70\%).\cite{III-1-2-ref27}

2. \textbf{Niveau 2 (Consolidation) :} Les 1 000 mots suivants (couverture \~85-90\%).

3. \textbf{Niveau 3 (Autonomie) :} Les 1 500 mots suivants (total \~3 000), spécifiquement
sélectionnés pour débloquer le seuil de 98\% sur les corpus cibles (NT, Platon,
Xénophon).\cite{III-1-2-ref4}



\subsection{Création et adaptation de textes "passerelles"}

Puisque le marché éditorial est déficitaire, la refondation doit inclure la production ou
l'adaptation de textes.

\begin{itemize}\item \textbf{Technique du \og Tiering\fg{} :} Réécrire des textes classiques en limitant le vocabulaire aux tranches définies ci-dessus. Par exemple, une version de l'\textit{Apologie de Socrate} limitée aux 1 500 mots les plus fréquents, puis une version \og enrichie\fg{} introduisant progressivement les mots rares.\item \textbf{Graded Readers Numériques :} Utiliser des plateformes (comme Biblingo ou des outils basés sur Alpheios) qui permettent de filtrer les textes ou de gloser automatiquement les mots hors-fréquence, garantissant que l'élève est toujours exposé à un ratio de 95-98\% de connu.\cite{III-1-2-ref52}\end{itemize}

\subsection{Pédagogie de l'inférence et de la morphologie}

L'enseignement explicite des stratégies de lecture doit être intégré :

\begin{itemize}\item \textbf{Entraînement à l'Inférence :} Ne pas laisser l'élève deviner au hasard. Enseigner comment utiliser le contexte syntagmatique et les indices morphologiques pour déduire le sens.\cite{III-1-2-ref54}\item \textbf{Routine de Parsing :} Automatiser l'analyse morphologique par des exercices de reconnaissance rapide (drills) pour réduire la latence de traitement des formes fléchies opaques.\cite{III-1-2-ref53}\end{itemize}

\subsection{Conclusion et recommandations opérationnelles}

L'analyse approfondie du point III.1.\cite{III-1-2-ref2} confirme que l'objectif de \og 95-98\% de
mots connus\fg{} est scientifiquement fondé et indispensable pour une refondation réussie de
l'enseignement du grec ancien.

1. \textbf{Validité Scientifique :} La recherche (Hu \& Nation, Laufer) démontre sans équivoque que
le seuil de \textbf{98\%} est le point de bascule vers la lecture fluide autonome. Le seuil de 95\%
doit être considéré comme une étape transitoire de lecture assistée, non comme une finalité.

2. \textbf{Dimensionnement de l'Effort :} Pour atteindre ce seuil sur des textes majeurs (Platon,
NT), le vocabulaire actif doit atteindre environ \textbf{3 100 à 3 600 lemmes}, soit le double des
standards actuels de fin de cycle secondaire.

3. \textbf{Contrainte Morphologique :} L'opacité de la morphologie grecque exige que cet
apprentissage lexical soit couplé à une conscience morphologique aiguë pour garantir la
reconnaissance des formes en contexte.

4. \textbf{Stratégie de Refondation :}

\begin{itemize}\item \textbf{Proscrire l'usage du dictionnaire} durant les phases de lecture extensive pour éviter la surcharge cognitive (split-attention).\item \textbf{Adopter des manuels à haute densité textuelle} (type \textit{Athenaze} Italien ou \textit{Alexandros}) et développer des corpus intermédiaires calibrés statistiquement.\item \textbf{Intégrer les outils numériques} (parsing à la demande) pour maintenir le flux de lecture lorsque le vocabulaire n'est pas encore acquis.\end{itemize}La lecture fluide en grec ancien est possible, mais elle ne peut être le fruit du hasard ou du seul
talent des élèves. Elle exige une architecture pédagogique rigoureuse, où chaque texte est une
marche calculée vers l'autonomie lexicale de 98\%.

\begin{table}[h!]\small\centering\begin{tabular}{| p{2cm} | p{2.2cm} | p{2.2cm} | p{2.2cm} |}\hline\textbf{Indicateur} & \textbf{Seuil Minimal (95\%)} & \textbf{Seuil Optimal (98\%)} & \textbf{Implication pour la Refondation} \\ \hline\textbf{Densité Mots Inconnus} & 1 mot sur 20 & 1 mot sur 50 & Privilégier les textes à faible densité d'inconnus pour la lecture cursive. \\ \hline\textbf{Type de Lecture} & Assistée / Intensive & Autonome / Extensive / Plaisir & Réserver le 95\% pour l'étude en classe ; viser 98\% pour le travail personnel. \\ \hline\textbf{Vocabulaire Requis (Platon/NT)} & \~1 750 lemmes & \~3 100 lemmes & Fixer un objectif lexical de fin de cycle à 3 000+ mots. \\ \hline\textbf{Rôle du Contexte} & Inférence difficile, risque d'erreur & Inférence fiable et automatique & Enseigner explicitement les stratégies d'inférence contextuelle. \\ \hline\textbf{Outil de Référence} & Dictionnaire nécessaire (Split-attention) & Aucun ou glose marginale & Supprimer le dictionnaire pour l'entraînement à la fluidité ; utiliser des gloses intégrées. \\ \hline\end{tabular}\end{table}